%%  Technical Chapters

\chapter{System Overview}

This project builds a system that uses EEG signals recorded from the scalp to do two
things. The first is to tell whether a person is imagining a movement of the left hand or
the right hand. The second is to find out the mental condition of the person at that
moment. Both outputs come from the same recording. The system takes the raw brain signals,
cleans them, pulls out useful information from them, and then uses a machine learning
model for the left versus right decision and a set of frequency band calculations for the
cognitive score. No audio is recorded at any stage.\\

\begin{figure}[h]
    \centering
    \includegraphics[width=0.60\textwidth]{BDeeg.png}
    \caption{EEG Based Cognitive Assessment Overview}
    \label{fig:system_overview}
\end{figure}

Figure~\ref{fig:system_overview} illustrates the overall architecture of the proposed
system. The first step is collecting the data. The person wears an EEG headset with
electrodes placed on the scalp at standard positions. During the recording session the
person is shown an instruction and is asked to imagine gripping an object with the left
hand or the right hand, while the EEG signals are being recorded. The task is only
imagined, so no actual movement is made and no sound is produced. To know which task was
being performed at each point in the recording, markers are placed in the recording
software at the start of every trial. These markers are what label the data, and they
replace the need for any audio reference.

The filtered EEG is then sent to the band analysis stage. Here the signal is divided into
the standard frequency bands, which are delta, theta, alpha, beta and gamma, to study the
brain activity. The bands matter for both halves of the system. For motor imagery, the
alpha and beta bands over the motor cortex are the ones that change when a movement is
imagined. For the cognitive score, the balance of power between the bands is what
indicates whether the person is attentive, under load, relaxed, stressed or tired.

The next step is feature extraction, which means taking out useful information from the
EEG signals. For the motor imagery path, the two motor cortex channels C3 and C4 are
filtered to the 8--30 Hz range, their power is computed, and the change in power with
respect to a baseline taken just before the task is calculated. Statistical measures such
as the mean, median, skewness and kurtosis of each channel are then worked out, along with
the ratio of power between C3 and C4 and the band power over 8--30 Hz. These values
together form a set of numbers that represent what the motor cortex was doing during that
trial. For the cognitive path, the power spectral density is computed and the relative
power in each band is obtained by dividing the power in that band by the total power.

The machine learning model is trained using the features extracted from the EEG data along
with the labels that say whether the trial was a left-hand or a right-hand imagination. The
model looks at the patterns in the features and tries to learn which pattern belongs to
which side. Once training is done, the model can take a new EEG trial and predict whether
the person imagined the left or the right hand. In parallel, the relative band powers are
combined into ratios to produce the cognitive outputs, which are attention, cognitive load,
relaxation, stress and fatigue.

The output of the system is therefore the predicted side, left or right, along with the
cognitive score. The predicted side shows which hand the person was imagining, and the
cognitive score shows how attentive, stressed or tired the person is at that moment. Both
outputs come from the same EEG recording setup, which keeps the system simple and practical
to use.

Overall this system takes EEG signals as input, processes them, and uses machine learning
and frequency band analysis to give two useful outputs. It works with a simple wearable
headset and needs nothing besides the EEG signal itself, which makes it a practical
approach for real world use.

%% new chapter

\chapter{System Hardware}
This chapter describes the hardware components used in the proposed EEG-based cognitive
assessment system. The components are used to record EEG signals, process the data in real
time, and display the output through the GUI. This chapter explains the use of the GELO BCI
headset, Raspberry Pi 5, display unit, and other supporting hardware used for brain signal
recording and cognitive analysis.

\section{EEG Recorder}
EEG recorders are used to record the electrical activity of the brain using electrodes
placed on the scalp. Different EEG systems are available based on factors like signal
quality, portability, comfort, number of channels, and communication. Depending on the
application, a suitable EEG recorder system is selected.

There are different types of EEG recorder system:

\subsection{Wet Electrode EEG System}

Wet electrode EEG recorder systems use conductive gel between the scalp and the electrode
to obtain better signal quality. These EEG recorder systems are mainly used for clinical
and medical purposes because they provide accurate and stable EEG recordings. In hospitals
and research laboratories, wet electrode EEG recorder systems commonly follow the standard
10--20 electrode placement system \cite{jasper1958tentwenty} for recording brain activity
from different regions of the scalp.

However, wet electrode EEG recorder systems are not suitable for the proposed project.
Before every recording session, conductive gel has to be applied properly, which increases
the setup time. The system also becomes less comfortable during long-duration usage. In
addition, regular cleaning and maintenance are required after recording. Since the proposed
project mainly focuses on wearable and real-time applications, wet electrode EEG recorder
systems are not very practical.

\subsection{Dry Electrode EEG System}

Dry electrode EEG recorder systems record EEG signals without using conductive gel. These
systems are easier to use and more comfortable compared to wet electrode systems. Dry
electrode EEG recorder systems are widely used in wearable Brain Computer Interface (BCI)
applications because they support faster setup and better portability.

Different types of dry electrode EEG recorder systems are available, such as EEG headbands
and EEG headphones. EEG headbands are lightweight and easy to wear, making them suitable
for continuous monitoring applications. EEG headphones provide portability and better
comfort. For this project, a wearable dry electrode EEG headset was preferred because it is
more comfortable, portable, and easy to use. Based on these requirements, the GELO BCI
headset was selected for recording EEG signals.

\subsection{GELO BCI Headset}
The GELO BCI Digital EEG Recorder Headphones shown in Figure~\ref{fig:gelo} serve as the
primary EEG signal acquisition unit of the proposed system. The device combines a wearable
headphone with integrated dry electrodes, enabling comfortable and non-invasive brain
signal recording without the need for conductive gel. It consists of 4 EEG channels, namely
A1, A2, C3 and C4, together with a reference electrode placed at Cz and two ground
electrodes. The placement of the electrodes is based on the standard 10--20 EEG electrode
system \cite{gelo2026manual}.

The channel layout is the main reason this device suits the project. C3 and C4 sit directly
over the left and right motor cortex, which are the two positions whose activity changes
when a hand movement is imagined, so they carry the information needed for the left versus
right decision. A1 and A2 sit at the ears and are used as the second terminal of each pair,
so that the system measures the difference A1--C3 and A2--C4 rather than an absolute
voltage. Taking a difference in this way cancels out noise that is common to both
electrodes. The same four channels also supply the frequency band powers used for the
cognitive score.

The headset operates at a sampling frequency of 250 Hz, which is high enough for the
0.5--35 Hz range of interest in this work. The device supports Bluetooth BLE 5.0
connectivity for wireless real-time data transmission to the companion computer, ensuring
portability and ease of use. The device also includes stereo speakers and a microphone as
part of its headphone function. These are not used in the present work, since no audio is
recorded, but they are available if an audio cue needs to be played to the subject during a
trial. These features make the GELO BCI headset well suited to the EEG-based motor imagery
classification and cognitive score assessment carried out in the proposed system.

\subsection{Features of GELO BCI Headset}

\begin{itemize}
    \item Audio headphone with integrated neuro-headset design
    \item 4 EEG channels: A1, A2, C3, and C4 with REF--CZ reference electrode
    \item Dry gilded PogoPin sensor material
    \item Sampling frequency of 250\,Hz
    \item Bluetooth Low Energy (BLE 5.0) communication interface
    \item Supports bipolar read mode
    \item Reproducible frequency range: 20--20000\,Hz
    \item Dual 40\,mm speakers
    \item Bluetooth profiles: HFP, HSP, A2DP, and VCRCP
\end{itemize}

\begin{figure}[!htbp]
    \centering
    \includegraphics[width=0.65\linewidth]{Pic1.jpeg}
    \caption{GELO BCI Headset}
    \label{fig:gelo}
\end{figure}

\newpage

\section{Raspberry Pi 5}

The Raspberry Pi 5 shown in Figure~\ref{fig:rpi5} serves as the main processing unit of the
proposed project. It is used to receive EEG signals from the GELO BCI headset through
Bluetooth and to perform feature extraction and machine learning classification. It also
handles cognitive score calculation and GUI-based visualization of EEG signals and
predicted outputs.

The board is powered by a Broadcom BCM2712 quad-core 64-bit ARM Cortex-A76 processor
clocked at 2.4\,GHz \cite{rpi5brief}, which provides sufficient computational performance for real-time EEG
analysis applications. The Raspberry Pi 5 supports interfaces such as GPIO, I2C, SPI, UART,
USB and PCIe, allowing easy integration with EEG devices, display modules and other
peripherals. It also supports dual micro-HDMI output for high-resolution display and
real-time monitoring.

In addition, the board supports dual-band Wi-Fi, Bluetooth 5.0/BLE, Gigabit Ethernet and
USB 3.0 connectivity, making wireless communication and peripheral interfacing easier.
These features make the Raspberry Pi 5 suitable for the proposed EEG-based motor imagery
classification and cognitive assessment system.

\begin{itemize}
    \item 2.4\,GHz quad-core ARM Cortex-A76 64-bit processor
    \item Available in 4\,GB or 8\,GB LPDDR4X RAM variants
    \item Dual-band 2.4\,GHz/5\,GHz Wi-Fi and Bluetooth 5.0 / BLE
    \item Gigabit Ethernet and USB 3.0 support
    \item Dual micro-HDMI ports with dual 4K display support
    \item 40-pin GPIO header for peripheral interfacing
    \item PCIe 2.0 x1 support for high-speed peripherals
    \item Real-time clock (RTC) and onboard power button
    \item Low power consumption suitable for continuous embedded operation
\end{itemize}

\begin{figure}[!htbp]
    \centering
    \includegraphics[width=0.65\linewidth]{pi.jpg}
    \caption{Raspberry Pi 5}
    \label{fig:rpi5}
\end{figure}

\newpage

\section{Aluminum Heat Sink Case with Dual Fans for Raspberry Pi 5}

The Aluminum Heat Sink Case with Dual Fans shown in Figure~\ref{fig:heatsink} is used to
provide efficient thermal management for the Raspberry Pi 5. Due to the computational load
involved in continuous EEG processing and classification, proper heat dissipation is
essential to maintain stable performance and prevent thermal throttling.

This case is constructed using CNC-machined aluminum, which acts as a passive heat sink.
Additionally, it includes two high-speed cooling fans that provide active airflow, ensuring
the temperature of the Raspberry Pi remains within safe operating limits. The case also
provides physical protection and structural rigidity, which is useful because the unit is
meant to be carried and used alongside a wearable headset rather than kept in a fixed
laboratory setup.

\subsection{Features of Aluminum Heat Sink}

\begin{itemize}
    \item Aluminum alloy body for superior passive heat dissipation
    \item Dual cooling fans for improved active airflow
    \item Reduced thermal throttling under heavy processing loads
    \item Protects Raspberry Pi from dust and physical damage
\end{itemize}

\begin{figure}[!htbp]
    \centering
    \includegraphics[width=0.55\linewidth]{sfs_03_2.jpg}
    \caption{Aluminum Heat Sink Case with Dual Fans for Raspberry Pi 5}
    \label{fig:heatsink}
\end{figure}

\section{Waveshare 7-Inch LCD Display}

The Waveshare HDMI Display \cite{waveshare7inch} shown in Figure~\ref{fig:lcd} is used to
provide the visual user interface for the proposed EEG-based cognitive assessment system. The display allows users
to interact with the graphical interface, visualize real-time EEG waveforms, follow the
left and right motor imagery cues during a recording session, and observe outputs such as
the predicted side and the cognitive score. A clear and responsive display is essential for
effective visualization of EEG signal variations, system status and real-time analysis
results.

The display supports HDMI input and is fully compatible with the Raspberry Pi 5, making it
suitable for embedded real-time GUI applications. Its capacitive touch support enables
smooth user interaction for controlling recording, calibration and signal visualization.
The compact and portable design makes it suitable for practical deployment, live
demonstrations and continuous monitoring applications in EEG-based human--machine
interaction systems.

\subsection{Features of Waveshare LCD Display}
\begin{itemize}
    \item 7-inch LCD screen with clear and sharp display output
    \item HDMI interface compatible with Raspberry Pi 5
    \item Supports Raspberry Pi OS and other operating systems
    \item Compact and lightweight design for embedded applications
    \item Enables real-time visual feedback and user interaction
\end{itemize}

\begin{figure}[!htbp]
    \centering
    \includegraphics[width=0.55\linewidth]{Exterior-Size.jpg}
    \caption{Waveshare 7-Inch LCD Display}
    \label{fig:lcd}
\end{figure}

\chapter{System Software}

This chapter explains the software setup used to develop the proposed EEG-based motor
imagery classification and cognitive assessment system. It includes the software
environment, the GUI software, the EEG signal collection process, and the libraries used
for signal processing, visualization, machine learning and cognitive score calculation. The
system combines real-time EEG data collection using the Capsule GUI with Python-based
preprocessing and classification methods to perform EEG signal analysis in a reliable way.

\section{Software Requirements}

The proposed system uses the Capsule GUI software for real-time EEG signal acquisition,
electrode impedance checking, calibration, signal display and marker placement. All
recording and processing reported in this work is carried out on a Windows 11 computer. The
Raspberry Pi 5 described in the previous chapter is the intended deployment platform, so
that the complete system can later run as a standalone unit, but the results presented in
this report were obtained on the Windows setup. The software requirements below list the
tools, programming environments and libraries needed to develop and run the system.

\begin{itemize}
    \item \textbf{Operating System:} Windows 11 for recording and analysis; Raspberry Pi OS
          for the planned embedded deployment
    \item \textbf{Programming Language:} Python 3.x
    \item \textbf{GUI Software:} Capsule software for EEG device connection, impedance
          check, calibration, recording and marker placement
    \item \textbf{EEG Device Support:} GELO BCI Bluetooth BLE 5.0 communication
    \item \textbf{Additional Tool:} MATLAB, used for the event-related desynchronisation
          and synchronisation analysis of the reference dataset and for generating the
          corresponding plots
    \item \textbf{Required Libraries and Tools:}
    \begin{itemize}
        \item NumPy -- Used for numerical computation on arrays of EEG samples, adding up
              spectrum bins into bands, and computing ratios
        \item SciPy -- Used for Butterworth filtering, FFT, Welch power spectral density
              estimation and spectral feature extraction from the raw EEG
        \item Pandas -- Used for structured EEG dataset management, trial labelling and
              handling of the converted CSV files
        \item h5py -- Used to read the recorded \texttt{.h5} session files so they can be
              converted to CSV for further analysis
        \item MNE-Python \cite{gramfort2013mne} -- Used for EEG preprocessing, channel
              handling and segmenting the continuous recording into epochs around each
              marker
        \item Matplotlib -- Used for the power spectral density plots, ERD curves,
              frequency band bars and the figures presented in this report
        \item PyQtGraph -- Used for the live EEG trace, where the plot has to refresh
              continuously as new samples arrive
        \item Scikit-learn \cite{pedregosa2011sklearn} -- Used for Support Vector Machine
              based classification of left versus right motor imagery, along with feature
              scaling and accuracy evaluation
        \item CustomTkinter -- Used to build the application interface, including the
              sidebar, cards, buttons and dialogs
        \item PyTorch (optional) -- Kept as an option for deep learning models in future
              work
    \end{itemize}
\end{itemize}

Matplotlib and PyQtGraph are both used because they serve different purposes. Matplotlib
produces the static figures used for analysis and reporting, while PyQtGraph is faster at
redrawing and is therefore used for the live signal display during a recording session.

These requirements ensure reliable functioning of the EEG acquisition, visualization,
classification and cognitive scoring parts of the system.

\vspace{0.5cm}

\section{Flowchart}
The system shown in Figure~\ref{fig:Flowchart} starts with the person putting on the EEG
headset and the headset being paired with the computer over Bluetooth. The resistance of
each of the four electrodes is then checked, since a poor contact makes the rest of the
recording useless. Two short calibration steps follow. In the first, the subject keeps the
eyes closed for 30 seconds so that the individual alpha peak frequency can be found, as
alpha activity is strongest with the eyes closed. In the second, the subject looks at a
fixed point for 30 seconds so that the median power in each frequency band can be stored
as a baseline. Later readings are compared against this baseline rather than against fixed
values, which makes the scores specific to that person.

\begin{figure}[!htbp]
    \centering
    \includegraphics[width=1\linewidth]{flowchart_1.png}
    \caption{Flowchart of EEG based Cognitive Assessment Model}
    \label{fig:Flowchart}
\end{figure}
\newpage

\FloatBarrier

Once calibration is complete, the recording begins and the subject performs the motor
imagery trials. A marker is placed at the start of each trial so that every segment of the
recording is tied to a known label, either left-hand or right-hand imagination. The session
is then saved in \texttt{.h5} format and converted to CSV using a Python script for further
analysis. No audio is recorded at any point.

If the signal is too noisy or unclear, the recording is done again from the beginning. Only
when the signal quality is good enough does the system move forward. After passing the
quality check, the pipeline splits into two separate paths that run independently.

The first path handles cognitive assessment. Here the EEG signal is broken down into the
five frequency bands, which are delta, theta, alpha, beta and gamma. The power of each band
is calculated and divided by the total power to give its relative power, which removes the
effect of the overall signal level changing as electrode contact improves or worsens. These
relative powers are then combined into ratios to give the cognitive outputs. Attention,
cognitive load and relaxation are read from the device SDK, while stress is obtained from
the ratio of relative high-beta power to relative alpha power, and fatigue from the ratio
of combined theta and alpha power to beta power. Each ratio is compared against the
person's own reference value and mapped onto a score between 0 and 100, so the result tells
how mentally active or tired the person is at that moment.

The second path handles motor imagery classification. In this path the C3 and C4 channels
are taken out and filtered to the 8--30 Hz band, a fixed window of 4 seconds is cut from
the start of each imagination, and the power and the change in power relative to the
baseline are computed. Features such as the mean, median, skewness and kurtosis of each
channel, the ratio of power between C3 and C4, and the band power over 8--30 Hz are then
extracted. These features are used to train a Support Vector Machine. After training, the
model accuracy is checked. If the accuracy is not satisfactory, the process goes back to
the feature extraction stage, the feature set or the kernel is adjusted, and the model is
retrained. This loop continues until the accuracy reaches an acceptable level. Once the
accuracy is good, the model is used to predict left or right from new EEG recordings.

Both paths finish at the end node, giving two outputs from a single EEG recording setup,
which are the cognitive score and the predicted side.